% 02-co-so-ly-thuyet.tex — Chương 2: Cơ sở lý thuyết và tổng quan
\chapter{Cơ sở lý thuyết và tổng quan}
\label{ch:coso}

Chương này hệ thống hóa các kiến thức nền tảng cần thiết cho đề tài. Trước hết, mục
\ref{sec:chatbot} giới thiệu khái niệm chatbot và trợ lý ảo. Mục \ref{sec:llm} trình bày
mô hình ngôn ngữ lớn và kiến trúc Transformer. Mục \ref{sec:rag} và mục \ref{sec:ft} lần
lượt đi sâu vào hai kỹ thuật trọng tâm là sinh tăng cường truy hồi (\rag) và tinh chỉnh
(\ft). Mục \ref{sec:hybrid} so sánh hai hướng tiếp cận và đề xuất kiến trúc lai. Mục
\ref{sec:metrics} trình bày các thước đo đánh giá, và mục \ref{sec:infra} giới thiệu các
nền tảng triển khai thực tế.

\section{Chatbot và trợ lý ảo}
\label{sec:chatbot}

\subsection{Khái niệm và phân loại}

Chatbot là một chương trình máy tính được thiết kế để tương tác với con người thông qua
hội thoại ngôn ngữ tự nhiên. Có thể phân loại chatbot theo nhiều tiêu chí. Theo mức độ
phức tạp của cơ chế xử lý, chatbot được chia thành ba thế hệ:

\begin{itemize}
  \item \textbf{Chatbot dựa trên luật (rule-based):} hoạt động theo các kịch bản và mẫu
  câu lệnh được định nghĩa trước. Ưu điểm là dễ kiểm soát và tiên đoán, nhưng cứng nhắc,
  khó mở rộng và dễ thất bại trước các cách diễn đạt ngoài kịch bản.
  \item \textbf{Chatbot dựa trên học máy (retrieval-based / intent-based):} sử dụng các
  mô hình phân loại ý định và trích xuất thực thể để ánh xạ câu hỏi người dùng vào một tập
  câu trả lời định sẵn. Linh hoạt hơn rule-based nhưng vẫn phụ thuộc vào tập ý định hữu hạn.
  \item \textbf{Chatbot sinh văn bản (generative) dựa trên mô hình ngôn ngữ lớn:} sinh câu
  trả lời trực tiếp từ mô hình ngôn ngữ, có khả năng diễn đạt tự nhiên và xử lý câu hỏi mở.
  Đây là hướng tiếp cận mà đề tài lựa chọn.
\end{itemize}

\subsection{Trợ lý ảo và kiến trúc hội thoại}

Khái niệm ``trợ lý ảo'' (virtual assistant) thường được hiểu rộng hơn chatbot: không chỉ
trả lời câu hỏi mà còn có khả năng thực hiện tác vụ, gọi công cụ, truy cập hệ thống nghiệp
vụ và duy trì ngữ cảnh qua nhiều lượt hội thoại. Một trợ lý ảo hiện đại thường gồm các
thành phần: (i) bộ hiểu ngôn ngữ và quản lý hội thoại; (ii) bộ truy xuất tri thức; (iii)
bộ gọi công cụ (tool calling); (iv) bộ sinh câu trả lời; và (v) lớp quan sát, ghi log để
giám sát chất lượng.

Trong các hệ thống lớn, kiến trúc \emph{đa tác tử} (multi-agent) ngày càng phổ biến: một
tác tử điều phối (orchestrator/host) tiếp nhận yêu cầu người dùng và định tuyến đến các
tác tử chuyên biệt (specialist agents), mỗi tác tử phụ trách một miền nghiệp vụ hẹp. Cách
tổ chức này giúp hệ thống dễ mở rộng, dễ kiểm thử và giảm tải nhận thức cho từng thành
phần. Để các tác tử có thể giao tiếp và chia sẻ công cụ một cách chuẩn hóa, các giao thức
như Model Context Protocol (MCP) và Agent-to-Agent (A2A) được sử dụng (trình bày ở mục
\ref{sec:infra}).

\section{Mô hình ngôn ngữ lớn}
\label{sec:llm}

\subsection{Kiến trúc Transformer}

Hầu hết các mô hình ngôn ngữ lớn hiện đại đều dựa trên kiến trúc Transformer
\cite{vaswani2017attention}. Khác với mạng hồi tiếp (RNN) xử lý tuần tự từng token,
Transformer xử lý toàn bộ chuỗi song song nhờ cơ chế \emph{tự chú ý} (self-attention),
cho phép mô hình nắm bắt quan hệ giữa các token ở xa nhau một cách hiệu quả và song song
hóa tốt trên GPU.

Cơ chế tự chú ý được mô tả bằng công thức \emph{scaled dot-product attention}:

\begin{equation}
\mathrm{Attention}(Q, K, V) = \mathrm{softmax}\!\left(\frac{QK^{\top}}{\sqrt{d_k}}\right) V,
\end{equation}

trong đó $Q$, $K$, $V$ lần lượt là các ma trận truy vấn (query), khóa (key) và giá trị
(value) được chiếu tuyến tính từ biểu diễn đầu vào, và $d_k$ là số chiều của vector khóa.
Cơ chế \emph{đa đầu} (multi-head attention) chạy song song nhiều phép chú ý trên các không
gian con khác nhau rồi ghép kết quả, giúp mô hình học nhiều loại quan hệ đồng thời. Mỗi
khối Transformer còn gồm mạng truyền thẳng theo vị trí (position-wise feed-forward
network), chuẩn hóa lớp (layer normalization) và kết nối tắt (residual connection).
Thông tin vị trí của token được bổ sung qua positional encoding.

\begin{figure}[htbp]
\centering
\begin{tikzpicture}[
  node distance=6mm,
  box/.style={draw=vnpostgreen,fill=boxblue,rounded corners,minimum width=4.2cm,minimum height=8mm,align=center,font=\small},
  att/.style={draw=vnpostdark,fill=white,rounded corners,minimum width=4.2cm,minimum height=8mm,align=center,font=\small},
  arrow/.style={-{Latex[length=2mm]},thick,vnpostdark}
]
\node[box] (input) {Token + Positional Encoding};
\node[att,above=of input] (mha) {Multi-Head Self-Attention};
\node[box,above=of mha] (add1) {Add \& Layer Norm};
\node[att,above=of add1] (ffn) {Feed-Forward Network};
\node[box,above=of ffn] (add2) {Add \& Layer Norm};
\node[box,above=of add2] (out) {Biểu diễn đầu ra};
\draw[arrow] (input) -- (mha);
\draw[arrow] (mha) -- (add1);
\draw[arrow] (add1) -- (ffn);
\draw[arrow] (ffn) -- (add2);
\draw[arrow] (add2) -- (out);
\draw[arrow] ($(input.west)+(-1.2,0)$) -- (input.west) node[midway,above,font=\scriptsize]{$x$};
\draw[arrow] ($(add2.east)+(1.2,0)$) -- (add2.east);
\end{tikzpicture}
\caption{Sơ đồ khối một khối Transformer (một lớp) với self-attention và feed-forward.}
\label{fig:transformer}
\end{figure}

\subsection{Tokenization và biểu diễn từ}

Trước khi đưa vào mô hình, văn bản được chia thành các \emph{token} bằng bộ tokenizer
(thường theo thuật toán BPE \cite{sennrich2016bpe} hoặc SentencePiece). Mỗi token được ánh
xạ thành một vector
nhúng; các vector này được cập nhật qua nhiều lớp Transformer để tạo ra biểu diễn ngữ cảnh
hóa. Chất lượng tokenizer ảnh hưởng trực tiếp đến khả năng xử lý tiếng Việt, đặc biệt với
các từ ghép và dấu thanh.

\subsection{Quá trình huấn luyện}

Một mô hình ngôn ngữ lớn thường trải qua các giai đoạn:

\begin{enumerate}
  \item \textbf{Tiền huấn luyện (pre-training):} huấn luyện mô hình dự đoán token tiếp theo
  trên khối lượng văn bản khổng lồ \cite{zhao2023survey}. Giai đoạn này giúp mô hình học
  ngữ pháp, kiến thức phổ
  quát và khả năng suy luận ngôn ngữ.
  \item \textbf{Tinh chỉnh có giám sát (instruction tuning/SFT):} tiếp tục huấn luyện trên
  các cặp hướng dẫn -- đáp án để mô hình biết cách ``trả lời'' thay vì chỉ ``hoàn thành
  văn bản'' \cite{chung2022instruction}.
  \item \textbf{Căn chỉnh theo phản hồi con người (RLHF) hoặc tối ưu theo sở thích (DPO):}
  điều chỉnh mô hình theo sở thích của con người về tính hữu ích, trung thực và an toàn
  \cite{ouyang2022instructgpt,rafailov2023dpo}.
\end{enumerate}

\subsection{Sinh văn bản và nhắc lệnh}

Khi suy luận, mô hình sinh token lần lượt theo một chiến lược giải mã. Nhiệt độ
(temperature) điều chỉnh độ ngẫu nhiên; top-$k$ và top-$p$ (nucleus sampling) giới hạn
tập token được chọn, giúp cân bằng giữa tính đa dạng và tính ổn định. Chất lượng đầu ra
phụ thuộc rất lớn vào \emph{nhắc lệnh} (prompt): kỹ thuật zero-shot, few-shot (đưa ví dụ)
và chuỗi suy luận (Chain-of-Thought -- CoT) \cite{wei2022cot} giúp định hướng mô hình giải
quyết tác vụ.

\section{Sinh tăng cường truy hồi (RAG)}
\label{sec:rag}

\subsection{Khái niệm và động lực}

Sinh tăng cường truy hồi \cite{lewis2020rag} là kỹ thuật kết hợp một bộ truy hồi
(retriever) với một mô hình sinh (generator). Thay vì buộc mô hình phải ``nhớ'' toàn bộ
tri thức trong tham số, \rag{} truy hồi các đoạn tài liệu liên quan từ kho tri thức ngoài
tại thời điểm suy luận, rồi cung cấp chúng làm ngữ cảnh để mô hình sinh câu trả lời. Cách
tiếp cận này mang lại ba lợi ích chính: (i) cập nhật tri thức mà không cần huấn luyện lại;
(ii) giảm hiện tượng ảo giác nhờ ràng buộc vào bằng chứng; và (iii) cho phép trích dẫn
nguồn, tăng tính kiểm chứng.

\subsection{Kiến trúc tổng quát}

Một hệ \rag{} điển hình gồm hai pha: \emph{lập chỉ mục} (indexing) và \emph{suy luận}
(inference). Ở pha lập chỉ mục, tài liệu được chuẩn hóa, chia đoạn, nhúng thành vector và
lưu vào cơ sở dữ liệu vector. Ở pha suy luận, câu hỏi của người dùng được nhúng, bộ truy
hồi tìm các đoạn liên quan, chúng được xếp hạng lại và ghép thành ngữ cảnh, cuối cùng mô
hình sinh câu trả lời.

\begin{figure}[htbp]
\centering
\begin{tikzpicture}[
  node distance=8mm,
  proc/.style={draw=vnpostgreen,fill=boxblue,rounded corners,minimum width=2.6cm,minimum height=9mm,align=center,font=\small},
  data/.style={draw=vnpostdark,fill=white,cylinder,shape border rotate=90,minimum width=2cm,minimum height=1.2cm,align=center,font=\small,aspect=0.35},
  arrow/.style={-{Latex[length=2mm]},thick,vnpostdark}
]
\node[proc] (q) {Câu hỏi};
\node[proc,right=of q] (ret) {Bộ truy hồi\\thưa + dày};
\node[data,right=of ret] (db) {Kho tri thức\\vector};
\node[proc,right=of db] (rank) {Xếp hạng lại};
\node[proc,right=of rank] (llm) {LLM sinh\\câu trả lời};
\node[proc,below=10mm of llm] (ans) {Câu trả lời\\có trích nguồn};
\draw[arrow] (q) -- (ret);
\draw[arrow] (ret) -- (db);
\draw[arrow] (db) -- (rank);
\draw[arrow] (rank) -- (llm);
\draw[arrow] (llm) -- (ans);
\draw[arrow] (db.south) to[out=-90,in=180] (rank.south);
\end{tikzpicture}
\caption{Kiến trúc tổng quát của một hệ thống RAG.}
\label{fig:rag-general}
\end{figure}

\subsection{Chia đoạn và tiền xử lý tài liệu}

Chất lượng của \rag{} phụ thuộc trước hết vào cách tài liệu được chia thành các \emph{đoạn}
(chunk). Đoạn quá lớn làm loãng ngữ cảnh và tăng chi phí; đoạn quá nhỏ làm mất mạch thông
tin. Với tài liệu nghiệp vụ chứa nhiều bảng biểu và tiêu đề, chiến lược chia đoạn thường
dùng là \emph{nhận biết cấu trúc} (header-aware): cắt theo tiêu đề mục, giữ nguyên bảng,
đặt độ dài tối đa và độ chồng lấn (overlap) hợp lý để bảo toàn ngữ cảnh giữa các đoạn liền
kề. Với tài liệu scan hoặc ảnh, cần bước nhận dạng ký tự quang học (OCR) trước khi chia
đoạn.

\begin{figure}[htbp]
\centering
\begin{tikzpicture}[
  node distance=7mm,
  proc/.style={draw=vnpostgreen,fill=boxblue,rounded corners,minimum width=2.3cm,minimum height=9mm,align=center,font=\small},
  arrow/.style={-{Latex[length=2mm]},thick,vnpostdark}
]
\node[proc] (doc) {Tài liệu\\gốc};
\node[proc,right=of doc] (ocr) {OCR /\\chuẩn hóa};
\node[proc,right=of ocr] (chunk) {Chia đoạn\\+ metadata};
\node[proc,right=of chunk] (emb) {Nhúng\\vector};
\node[proc,right=of emb] (store) {Lưu vào\\Qdrant};
\draw[arrow] (doc) -- (ocr);
\draw[arrow] (ocr) -- (chunk);
\draw[arrow] (chunk) -- (emb);
\draw[arrow] (emb) -- (store);
\end{tikzpicture}
\caption{Pipeline lập chỉ mục tài liệu trong hệ thống RAG.}
\label{fig:rag-index}
\end{figure}

\subsection{Nhúng vector và cơ sở dữ liệu vector}

Mô hình nhúng (embedding) ánh xạ đoạn văn bản thành một vector số thực trong không gian
nhiều chiều, sao cho các đoạn có ngữ nghĩa gần nhau nằm gần nhau. Độ tương đồng thường được
đo bằng cosine:

\begin{equation}
\mathrm{sim}(\mathbf{u}, \mathbf{v}) = \frac{\mathbf{u}\cdot\mathbf{v}}{\lVert\mathbf{u}\rVert\,\lVert\mathbf{v}\rVert}.
\end{equation}

Cơ sở dữ liệu vector (vector database) lưu trữ và tìm kiếm láng giềng gần nhất trên hàng
triệu vector một cách hiệu quả nhờ các thuật toán chỉ mục xấp xỉ \cite{malkov2020hnsw}.
Trong đề tài này, Qdrant \cite{qdrant2024} được sử dụng với khả năng lưu song
song vector dày (dense) và vector thưa (sparse) trong cùng một collection.

\subsection{Truy hồi thưa, truy hồi dày và truy hồi lai}

\emph{Truy hồi thưa} (sparse retrieval) biểu diễn văn bản bằng vector thưa dựa trên tần
suất từ, tiêu biểu là BM25 \cite{robertson2009bm25}. Điểm BM25 giữa câu truy vấn $q$ và
tài liệu $d$ được tính:

\begin{equation}
\mathrm{BM25}(q,d) = \sum_{t \in q} \mathrm{IDF}(t)\cdot
\frac{f(t,d)\,(k_1+1)}{f(t,d) + k_1\left(1-b+b\,\frac{|d|}{\mathrm{avgdl}}\right)},
\end{equation}

trong đó $f(t,d)$ là tần suất từ $t$ trong $d$, $|d|$ là độ dài tài liệu, $\mathrm{avgdl}$
là độ dài trung bình, và $k_1$, $b$ là các tham số điều chỉnh. BM25 mạnh ở việc khớp từ
khóa chính xác, đặc biệt với số hiệu văn bản, mã dịch vụ hoặc thuật ngữ hiếm.

\emph{Truy hồi dày} (dense retrieval) \cite{karpukhin2020dpr} sử dụng mô hình nhúng để
biểu diễn câu hỏi và tài liệu trong cùng không gian ngữ nghĩa, nhờ đó bắt được các trường
hợp diễn đạt đồng nghĩa mà BM25 bỏ sót. Hai phương pháp có ưu nhược điểm bổ trợ cho nhau.
\emph{Truy hồi lai} (hybrid retrieval) kết hợp cả hai bằng cách lấy hợp hoặc kết hợp điểm
số, thường cho kết quả tốt hơn từng phương pháp riêng lẻ.

\subsection{Xếp hạng lại}

Các bộ truy hồi giai đoạn đầu (first-stage retrieval) thường nhanh nhưng kém chính xác do
phải quét trên toàn kho. Bộ \emph{xếp hạng lại} (reranker) tinh chỉnh thứ tự ứng viên bằng
một mô hình mạnh hơn -- thường là mô hình cross-encoder đọc đồng thời câu hỏi và đoạn văn
để cho điểm tương quan. Quy trình điển hình: lấy một tập ứng viên lớn ở giai đoạn một (ví
dụ top-20 mỗi nhánh), hợp chúng lại, rồi xếp hạng lại toàn bộ và giữ top-$k$ nhỏ (ví dụ
top-5) làm ngữ cảnh cho mô hình sinh.

\begin{figure}[htbp]
\centering
\begin{tikzpicture}[
  node distance=7mm,
  proc/.style={draw=vnpostgreen,fill=boxblue,rounded corners,minimum width=2.4cm,minimum height=9mm,align=center,font=\small},
  arrow/.style={-{Latex[length=2mm]},thick,vnpostdark}
]
\node[proc] (q) {Câu hỏi};
\node[proc,above right=8mm and 14mm of q] (bm) {Truy hồi thưa\\BM25 (top-20)};
\node[proc,below right=8mm and 14mm of q] (de) {Truy hồi dày\\(top-20)};
\node[proc,right=32mm of q] (union) {Hợp ứng viên};
\node[proc,right=of union] (rr) {Xếp hạng lại\\cross-encoder};
\node[proc,right=of rr] (topk) {Top-$k$ ngữ cảnh};
\draw[arrow] (q) -- (bm);
\draw[arrow] (q) -- (de);
\draw[arrow] (bm) -- (union);
\draw[arrow] (de) -- (union);
\draw[arrow] (union) -- (rr);
\draw[arrow] (rr) -- (topk);
\end{tikzpicture}
\caption{Pipeline truy hồi lai kết hợp BM25, truy hồi dày và xếp hạng lại.}
\label{fig:rag-hybrid}
\end{figure}

\subsection{Các biến thể của RAG}

Từ kiến trúc cơ bản, \rag{} có nhiều biến thể \cite{gao2023ragsurvey}. \emph{Naive RAG} là dạng đơn giản nhất với
truy hồi một lần. \emph{Advanced RAG} bổ sung tiền xử lý truy vấn, hậu xử lý kết quả và
xếp hạng lại. \emph{Modular RAG} tổ chức thành các mô-đun linh hoạt có thể thay thế. Ngoài
ra, \emph{GraphRAG} khai thác đồ thị tri thức để truy hồi theo quan hệ giữa các thực thể,
và \emph{agentic RAG} cho phép tác tử quyết định truy hồi nhiều vòng, cải viết truy vấn
hoặc gọi công cụ khác khi cần. Đề tài này áp dụng advanced/modular RAG với truy hồi lai và
xếp hạng lại, đồng thời định hướng mở rộng lên agentic và GraphRAG trong tương lai.

\section{Tinh chỉnh mô hình ngôn ngữ lớn (Fine-tuning)}
\label{sec:ft}

\subsection{Khái niệm và các mức độ}

Tinh chỉnh (\ft) là quá trình tiếp tục huấn luyện một mô hình đã được tiền huấn luyện trên
một tập dữ liệu chuyên biệt, nhằm điều chỉnh mô hình phù hợp với miền hoặc tác vụ cụ thể.
Có thể chia thành hai nhóm:

\begin{itemize}
  \item \textbf{Tinh chỉnh toàn bộ (full fine-tuning):} cập nhật toàn bộ tham số mô hình.
  Hiệu quả cao nhưng đòi hỏi tài nguyên tính toán và bộ nhớ rất lớn, đồng thời cần lưu
  toàn bộ bản sao mô hình cho mỗi tác vụ.
  \item \textbf{Tinh chỉnh hiệu quả tham số (PEFT):} chỉ cập nhật một phần nhỏ tham số,
  giữ nguyên phần lớn mô hình gốc. Tiêu biểu là LoRA và QLoRA.
\end{itemize}

\subsection{Tinh chỉnh có giám sát và căn chỉnh}

\emph{Tinh chỉnh có giám sát} (Supervised Fine-Tuning -- SFT) huấn luyện mô hình trên các
cặp hướng dẫn -- đáp án chất lượng cao. Để mô hình trả lời đúng ``giọng'' nghiệp vụ, dữ
liệu SFT cần được thiết kế cẩn thận cả về nội dung lẫn định dạng. Sau SFT, có thể áp dụng
RLHF hoặc DPO \cite{rafailov2023dpo} để căn chỉnh theo sở thích con người, giúp câu trả
lời hữu ích và an toàn hơn.

\subsection{LoRA}

LoRA (Low-Rank Adaptation) \cite{hu2021lora} giả định rằng thay đổi tham số khi thích ứng
một tác vụ có ``hạng thấp'' (low intrinsic rank). Thay vì cập nhật toàn bộ ma trận trọng
số $W_0 \in \mathbb{R}^{d \times k}$, LoRA giữ nguyên $W_0$ và thêm một phần hiệu chỉnh
phân rã thành tích hai ma trận hạng thấp:

\begin{equation}
W = W_0 + \Delta W = W_0 + BA, \qquad B \in \mathbb{R}^{d \times r},\ A \in \mathbb{R}^{r \times k},\ r \ll \min(d,k).
\end{equation}

Ma trận $A$ được khởi tạo ngẫu nhiên và $B$ được khởi tạo bằng 0 để lúc bắt đầu
$\Delta W = 0$, bảo toàn hành vi mô hình gốc. Số tham số huấn luyện giảm từ $d \times k$
xuống $r(d+k)$, thường chỉ còn dưới 1\% tổng tham số. Hệ số $\alpha$ và hạng $r$ quyết
định độ lớn của cập nhật; giá trị hiệu chỉnh thường được nhân với $\alpha/r$.

\begin{figure}[htbp]
\centering
\begin{tikzpicture}[
  node distance=8mm,
  mat/.style={draw=vnpostdark,fill=white,minimum width=1.2cm,minimum height=2.2cm,align=center,font=\small},
  small/.style={draw=vnpostgreen,fill=boxblue,minimum width=0.9cm,minimum height=1.4cm,align=center,font=\small},
  arrow/.style={-{Latex[length=2mm]},thick,vnpostdark}
]
\node[mat] (w0) {$W_0$\\ (đóng băng)};
\node[small,right=18mm of w0] (b) {$B$};
\node[small,right=4mm of b] (a) {$A$};
\node[mat,right=18mm of a] (w) {$W = W_0 + BA$};
\node[font=\small,below=2mm of b] {$d\times r$};
\node[font=\small,below=2mm of a] {$r\times k$};
\draw[arrow] (w0.east) -- (w.west);
\draw[arrow] (b.east) -- (a.west);
\draw[arrow] (a.east) -- (w.west);
\end{tikzpicture}
\caption{Minh họa LoRA: giữ nguyên $W_0$ và thêm cập nhật hạng thấp $BA$.}
\label{fig:lora}
\end{figure}

\subsection{QLoRA}

QLoRA \cite{dettmers2023qlora} mở rộng LoRA bằng cách lượng tử hóa mô hình gốc xuống 4 bit
(định dạng NF4), áp dụng lượng tử hóa kép (double quantization) và tối ưu hóa phân trang
(paged optimizers). Nhờ đó có thể tinh chỉnh các mô hình lớn trên phần cứng khiêm tốn hơn
đáng kể so với full fine-tuning, với mức suy giảm chất lượng nhỏ.

\subsection{Khi nào nên dùng fine-tuning}

Fine-tuning phù hợp khi cần điều chỉnh \emph{hành vi} của mô hình (phong cách, định dạng,
quy trình nghiệp vụ ổn định) hoặc khi tri thức miền mang tính ổn định và khó mô tả bằng
ngữ cảnh truy hồi. Ngược lại, khi tri thức thay đổi thường xuyên hoặc cần trích dẫn nguồn,
\rag{} là lựa chọn phù hợp hơn.

\section{So sánh RAG và fine-tuning, kiến trúc lai}
\label{sec:hybrid}

\rag{} và \ft{} giải quyết hai lớp vấn đề khác nhau và có thể bổ trợ cho nhau. Bảng
\ref{tab:rag-vs-ft} tóm tắt sự so sánh.

\begin{table}[htbp]
\centering
\caption{So sánh RAG và fine-tuning theo các tiêu chí triển khai.}
\label{tab:rag-vs-ft}
\small
\begin{tabular}{@{}p{3.6cm}p{5.4cm}p{5.4cm}@{}}
\toprule
\textbf{Tiêu chí} & \textbf{RAG} & \textbf{Fine-tuning} \\
\midrule
Nguồn tri thức & Kho tài liệu ngoài, truy hồi động & Tham số mô hình, cố định sau huấn luyện \\
Cập nhật tri thức & Chỉ cần cập nhật chỉ mục & Cần huấn luyện lại \\
Khả năng trích nguồn & Có (dựa trên đoạn truy hồi) & Hạn chế \\
Chi phí huấn luyện & Thấp (chủ yếu lập chỉ mục) & Cao hơn (SFT/LoRA) \\
Chi phí suy luận & Cao hơn (thêm truy hồi) & Thấp hơn (chỉ sinh) \\
Điều chỉnh phong cách/định dạng & Yếu & Mạnh \\
Giảm ảo giác & Tốt nếu truy hồi đúng & Trung bình \\
Phù hợp nhất với & Tri thức động, tra cứu & Hành vi miền ổn định \\
\bottomrule
\end{tabular}
\end{table}

Từ so sánh trên, đề tài đề xuất \emph{kiến trúc lai} kết hợp cả hai: sử dụng \rag{} làm
kênh cung cấp tri thức động và bằng chứng, đồng thời dùng mô hình đã tinh chỉnh để sinh
câu trả lời đúng phong cách, đúng định dạng nghiệp vụ và xử lý các tình huống ổn định.
Một bộ định tuyến (router) quyết định khi nào chỉ cần \rag, khi nào dùng mô hình tinh chỉnh,
và khi nào kết hợp. Chi tiết kiến trúc được trình bày ở chương \ref{ch:kientruc}.

\section{Các thước đo đánh giá}
\label{sec:metrics}

Đánh giá một hệ \rag{} cần tách thành hai khâu: chất lượng truy hồi và chất lượng sinh.

\subsection{Thước đo truy hồi}

\begin{itemize}
  \item \textbf{Recall@$k$:} tỉ lệ các tài liệu/đoạn liên quan được tìm thấy trong top-$k$.
  Đây là thước đo quan trọng nhất vì nếu tài liệu đúng không nằm trong ngữ cảnh thì mô hình
  không thể trả lời đúng.
  \item \textbf{Precision@$k$:} tỉ lệ đoạn liên quan trong top-$k$.
  \item \textbf{MRR (Mean Reciprocal Rank):} nghịch đảo thứ hạng của đoạn liên quan đầu tiên,
  trung bình trên các truy vấn. MRR cao nghĩa là đoạn đúng xuất hiện sớm.
  \item \textbf{nDCG:} thước đo có trọng số theo vị trí, phù hợp khi đánh giá chất lượng
  xếp hạng với mức độ liên quan khác nhau.
\end{itemize}

\subsection{Thước đo sinh văn bản}

Bộ thước đo RAGAS \cite{es2023ragas} đánh giá hệ \rag{} không cần đáp án tham chiếu đầy đủ:

\begin{itemize}
  \item \textbf{Faithfulness (độ trung thực):} mức độ câu trả lời được hỗ trợ bởi ngữ cảnh
  truy hồi, phản ánh mức độ ảo giác.
  \item \textbf{Answer Relevancy (độ liên quan câu trả lời):} mức độ câu trả lời đúng trọng
  tâm câu hỏi.
  \item \textbf{Context Precision:} tỉ lệ đoạn ngữ cảnh thực sự hữu ích.
  \item \textbf{Context Recall:} mức độ ngữ cảnh bao phủ đáp án tham chiếu.
\end{itemize}

Ngoài ra, \emph{LLM-as-judge} sử dụng một mô hình ngôn ngữ mạnh để cho điểm câu trả lời
theo thang điểm (ví dụ 0--10) dựa trên đáp án chuẩn, là phương pháp linh hoạt và tiết kiệm
cho miền hẹp. Các thước đo truyền thống như BLEU và ROUGE có thể dùng bổ sung nhưng không
phản ánh tốt tính đúng đắn ngữ nghĩa trong hỏi đáp mở.

\section{Nền tảng triển khai}
\label{sec:infra}

\subsection{Phục vụ suy luận với vLLM}

vLLM \cite{kwon2023vllm,vllm2024} là thư viện phục vụ suy luận LLM hiệu năng cao, nổi bật với cơ chế
PagedAttention quản lý bộ nhớ KV-cache theo trang, cho phép tăng thông lượng đáng kể nhờ
xử lý theo lô liên tục (continuous batching). Kết hợp với lượng tử hóa FP8, vLLM cho phép
phục vụ các mô hình hàng chục tỷ tham số trên GPU với chi phí hợp lý.

\subsection{Lượng tử hóa FP8}

FP8 là định dạng số dấu phẩy động 8 bit, giảm một nửa dung lượng bộ nhớ so với FP16 và
tăng tốc tính toán trên phần cứng hỗ trợ, với mức suy giảm chất lượng thường nhỏ. Đây là
lựa chọn phổ biến để phục vụ mô hình lớn trên hạ tầng giới hạn.

\subsection{Giao thức MCP và A2A}

\textbf{Model Context Protocol (MCP)} \cite{mcp2024} chuẩn hóa cách mô hình truy cập công cụ và nguồn dữ
liệu, giúp tách biệt logic công cụ khỏi tác tử. \textbf{Agent-to-Agent (A2A)} \cite{a2a2024} chuẩn hóa
giao tiếp giữa các tác tử, cho phép một tác tử điều phối gọi tác tử chuyên biệt và nhận
kết quả. Hai giao thức này tạo nên nền tảng cho kiến trúc đa tác tử của đề tài.

\subsection{Quan sát hệ thống và quản lý prompt}

OpenTelemetry (OTel) \cite{otel2024} cung cấp chuẩn thu thập vết (trace), chỉ số (metric) và log; dữ liệu
được đẩy tới các nền tảng như Langfuse \cite{langfuse2024} hoặc Grafana Tempo để giám sát luồng suy luận, độ
trễ và lượng token. MLflow Prompt Registry \cite{mlflow2024} cho phép quản lý phiên bản prompt tập trung, giúp
tách cấu hình khỏi mã nguồn và kiểm soát thay đổi prompt theo môi trường (ví dụ nhãn
\texttt{@production}).

\section*{Kết chương}
\addcontentsline{toc}{section}{Kết chương}
Chương này đã hệ thống hóa nền tảng lý thuyết về chatbot, LLM, \rag, fine-tuning, kiến trúc
lai, thước đo đánh giá và nền tảng triển khai. Đây là cơ sở để chương tiếp theo phân tích
bài toán cụ thể của VNPost và thiết kế kiến trúc hệ thống đề xuất.
